\documentclass[12pt, a4paper]{article}

\usepackage[french]{babel}
\usepackage[utf8]{inputenc}
\usepackage[T1]{fontenc}
\usepackage{tgadventor}
\renewcommand*\familydefault{\sfdefault}
\usepackage{amsmath, amssymb}
\usepackage{geometry}
\usepackage{enumitem}
\usepackage{multicol}
\usepackage{xcolor}
\usepackage{mdframed}
\usepackage{verbatim}
\usepackage{tabularx}

\geometry{margin=2.5cm}

\definecolor{vlmblue}{RGB}{25, 90, 160}
\definecolor{vlmgray}{RGB}{245, 245, 245}
\definecolor{vlmgreen}{RGB}{0, 130, 80}

% ── Titre de l'exercice dans un encadré ──
\newcommand{\titrexercice}[2]{%
  \begin{mdframed}[
    backgroundcolor=vlmgray,
    linecolor=black,
    linewidth=1.5pt,
    innerleftmargin=10pt,
    innerrightmargin=10pt,
    innertopmargin=6pt,
    innerbottommargin=6pt,
    skipabove=1em,
    skipbelow=0.5em
  ]
  \textbf{Exercice #1}%
  \ifx&#2&\else\hfill\textbf{#2}\fi
  \end{mdframed}
}

% ── Corps de l'exercice sans cadre ──
\newenvironment{exercice}[2]{%
  \titrexercice{#1}{#2}
  \begin{list}{}{%
    \setlength{\leftmargin}{0pt}
    \setlength{\rightmargin}{0pt}
    \setlength{\listparindent}{0pt}
    \setlength{\itemindent}{0pt}
    \setlength{\parsep}{0pt}
  }\item[]\noindent\ignorespaces
}{%
  \end{list}
  \vspace{1em}
}

% ── Solution ──
\ifdefined\AVECSOLUTIONS
  \newenvironment{solution}[1]{%
    \vspace{0.3em}
    \textbf{\textcolor{vlmgreen}{Solution #1}}
    \vspace{0.3em}\par
  }{%
    \vspace{1em}
  }
\else
  \newenvironment{solution}[1]{\comment}{\endcomment}
\fi

\pagestyle{empty}

\begin{document}
\begin{exercice}{EX-CH01-005}{Nombres parfaits}

\begin{mdframed}[backgroundcolor=vlmgray, innerleftmargin=10pt, innerrightmargin=10pt, innertopmargin=6pt, innerbottommargin=6pt, skipabove=-5pt, skipbelow=5pt]
Un \textbf{nombre parfait} est un nombre naturel non nul qui est égal à la somme de ses diviseurs stricts (ses diviseurs sauf lui-même).
\end{mdframed}


\bigskip
Parmi les nombres suivants, lesquels sont parfaits ?

\medskip
\begin{tabularx}{\linewidth}{*{5}{>{\centering\arraybackslash}X}}
6 & 28 & 496 & 875 & 8128\\
\end{tabularx}

\end{exercice}
\begin{solution}{EX-CH01-005}

On détermine les diviseurs stricts de chaque nombre, puis on vérifie si leur somme est égale au nombre.

\bigskip

\begin{tabular}{llll}
$D^*(6)$    & $= \{1, 2, 3\}$                                                    & $1+2+3 = 6$    & $\Rightarrow$ \textbf{parfait} $\checkmark$ \\[6pt]
$D^*(28)$   & $= \{1, 2, 4, 7, 14\}$                                             & $1+2+4+7+14 = 28$  & $\Rightarrow$ \textbf{parfait} $\checkmark$ \\[6pt]
$D^*(496)$  & $= \{1, 2, 4, 8, 16, 31, 62, 124, 248\}$                           & $1+2+4+8+16+31+62+124+248 = 496$  & $\Rightarrow$ \textbf{parfait} $\checkmark$ \\[6pt]
$D^*(875)$  & $= \{1, 5, 7, 25, 35, 125, 175\}$                                  & $1+5+7+25+35+125+175 = 373 \neq 875$ & $\Rightarrow$ non parfait $\times$ \\[6pt]
$D^*(8128)$ & $= \{1, 2, 4, 8, 16, 32, 64, 127, 254, 508, 1016, 2032, 4064\}$   & $\sum = 8128$ & $\Rightarrow$ \textbf{parfait} $\checkmark$ \\[6pt]
\end{tabular}

\bigskip
Les nombres parfaits parmi ceux proposés sont donc : $\mathbf{6}$, $\mathbf{28}$, $\mathbf{496}$ et $\mathbf{8128}$.

\end{solution}
\end{document}